\documentclass[10pt,a4paper]{amsart}
\usepackage[margin=19mm]{geometry}
\usepackage{amsmath,amssymb,mathtools}
\usepackage[scr=rsfs,cal=euler]{mathalfa}
\usepackage{xfrac}
\usepackage[unicode,hidelinks,bookmarks=false]{hyperref}
\newcommand{\ie}{i.e.\ }
\newcommand{\cf}{cf.\ }
\newcommand{\resp}{respectively\ }
\newcommand{\Subsection}{\subsection}
\providecommand{\nobreakdash}{\nobreak\hbox{-}}
\setlength{\emergencystretch}{2em}
\multlinegap0pt
\usepackage{amssymb}
\usepackage{caption}
\usepackage{mathtools}
\usepackage{enumerate}
\usepackage{xcolor}
\usepackage{tikz}
\usepackage{algorithm}\usepackage{algpseudocode}
\usepackage[normalem]{ulem}
\usepackage{multirow}
\usepackage{booktabs}
\usepackage{natbib}
\newtheorem{proposition}{Proposition}
\DeclareMathOperator{\GL}{GL}
\title{On the Coupled Cluster Doubles Truncation Variety of Four Electrons}
\author{Fabian M. Faulstich, Vincenzo Galgano, Elke Neuhaus, Irem Portakal}
\date{Source: arXiv:2602.16580v1 (CC BY 4.0)}
\begin{document}
\maketitle

\noindent\emph{Source and licence.} On the Coupled Cluster Doubles Truncation Variety of Four Electrons. Version arXiv:2602.16580v1 (CC BY 4.0), distributed by arXiv under the Creative Commons Attribution 4.0 International licence, http://creativecommons.org/licenses/by/4.0/, as recorded in the arXiv licence metadata for this exact version and shown on the abstract page https://arxiv.org/abs/2602.16580v1. Full licence notice: \texttt{licenses/arxiv-2602.16580-CC-BY-4.0.txt}.\par\smallskip

\noindent\textit{Editorial scope.} Counted unit: the article's proposition sym2wedge2 (source0.tex lines 889--891). The article's complete original proof of it is delivered with it (source0.tex lines 892--919, one \texttt{proof} environment of 28 lines). No mathematics is written, repaired or re-derived: the statement and the proof are the article's own lines, in the article's own notation, and no retained line is substituted or re-flowed.  Retained uncounted as the source-local context the proof needs: lines 831--834; lines 836--839.  Nothing was omitted from any retained span, so the package carries no omission record.  No retained line contains a citation, so the wrapper carries no bibliography and no bibliography entry.  No retained line contains a figure environment, an image-inclusion command or drawing code, so no image is delivered and the package declares no asset dependency.  Editorial formatting only, in addition: an AMS \texttt{amsart} preamble that restates the article's own declarations and macros used by the retained lines, namely the environments \texttt{proposition} on separate counters, together with the standard packages those lines need.  The retained environments print in this document as Proposition 1 in order of appearance, whereas the article prints the same items with the numbering of its own full document.  The complete native source of the article as distributed in the arXiv e-print for this exact version is bundled unchanged as \texttt{sources/194900/source0.tex}.
\begin{align}
{\textstyle \bigwedge^k}[3,1] & = [4-k,1^k] & \forall k=0,\ldots, 3 \label{eq:wedge[3,1]} \\
[3,1] \otimes [3,1] & = [4] \oplus [3,1] \oplus [2,2] \oplus [2,1,1] .\label{eq:[3,1]*[3,1]}
\end{align}

\begin{align}\label{eq:sym[3,1]}
S^2[3,1] & = [3,1]\otimes [3,1] - {\textstyle \bigwedge^2}[3,1] = [4] \oplus [3,1] \oplus [2,2] \nonumber \\
S^2[2,1,1] & = [4] \oplus [2,1,1] \oplus [2,2] .
\end{align}

\begin{proposition}\label{prop:sign-rep in sym2wedge2}
Consider the action of $\mathfrak{S}_4$ on $W = F \oplus G$ that fixes $F\simeq \mathbb C^{m-4}$ and permutes the basis vectors of $G\simeq \mathbb C^4$. Then, as $\mathfrak{S}_4$-representation, the linear space $V_{(0,2,0,\ldots)}^{A_{n-5}}$ contains no copy of the sign-representation $[1^4]$.
\end{proposition}

\begin{proof}
The $\GL(W)$-representation $V_{(0,2,0,\ldots)}^{A_{m-1}}$ is the kernel of the $\GL(W)$-equivariant map  
\[ \begin{matrix}
S^2(\textstyle{\bigwedge^2W}) & \rightarrow & \textstyle{\bigwedge^4W} \\
\alpha\cdot \beta & \mapsto & \alpha \wedge \beta
\end{matrix} \ .\] 
Therefore, to prove the hypothesis, it is enough to show that the module $S^2(\bigwedge^2W)$ contains as many copies of $[1^4]$ as $\bigwedge^4W$. Again, we decompose $W=\sum_{\ell=1}^{m-4}[4]_\ell \oplus [4]_0 \oplus [3,1]_0= \mathbb C^{m-3}\oplus [3,1]$. \\
First, we compute the $2$nd exterior power of $W$.
\begin{align*}
{\textstyle \bigwedge^2}W & \  = {\textstyle \bigwedge^2}\left( \mathbb C^{m-3} \oplus [3,1]\right)  = {\textstyle \bigwedge^2}\mathbb C^{m-3} \oplus \left(\mathbb C^{m-3}\wedge [3,1] \right) \oplus {\textstyle \bigwedge^2}[3,1] \\
& \stackrel{\eqref{eq:wedge[3,1]}}{=} {\textstyle \binom{m-3}{2}}[4] \oplus (m-3)[3,1] \oplus [2,1,1] .
\end{align*}
Second, applying the $2$nd symmetric power, we get
\begin{align*}
S^2({\textstyle \bigwedge^2}W) & \ = S^2\left( {\textstyle \binom{m-3}{2}}[4] \oplus (m-3)[3,1] \oplus [2,1,1] \right) \\
& \ \simeq {\textstyle \binom{m-3}{2}^2}S^2[4] \oplus {\textstyle \binom{m-3}{2}^2}(m-3)[4]\cdot [3,1] \oplus {\textstyle \binom{m-3}{2}^2}[4]\cdot [2,1,1] \oplus  (m-3)^2 S^2[3,1] \\
& \ \quad \oplus (m-3)[3,1]\cdot [2,1,1] \oplus S^2[2,1,1] \\
& \stackrel{\eqref{eq:sym[3,1]}}{=} {\textstyle \binom{m-3}{2}^2}[4] \oplus {\textstyle \binom{m-3}{2}^2}(m-3)[3,1] \oplus {\textstyle \binom{m-3}{2}^2}[2,1,1] \oplus (m-3)^2\big( [4] \oplus [3,1] \oplus [2,2] \big) \\
& \ \quad \oplus (m-3)[3,1]\cdot [2,1,1] \oplus \big( [4] \oplus [2,1,1] \oplus [2,2] \big) .
\end{align*}
Therefore, all copies of $[1^4]$ inside $S^2(\bigwedge^2W)$ lie inside the summands $(m-3)[3,1]\cdot [2,1,1]$. To decompose this product, we use that $[2,1,1]=[3,1]\otimes [1^4]$ and the decomposition of $[3,1]\otimes [3,1]$:
\begin{align*}
(m-3)[3,1]\cdot [2,1,1] & \ = (m-3) [3,1] \otimes \big([3,1] \otimes [1^4] \big) \\
& \stackrel{\eqref{eq:[3,1]*[3,1]}}{=} (m-3)\big( [4] \oplus [3,1] \oplus [2,2] \oplus [2,1,1]\big) \otimes [1^4]\\
& \ = (m-3)\big( [1^4] \oplus [2,1,1] \oplus [2,2] \oplus [3,1] \big) .
\end{align*}
We conclude that $S^2(\bigwedge^2W)$ has $m-3$ copies of $[1^4]$, all of them coming from the component $\bigwedge^4W$. In particular, $V_{(0,2,0,\ldots)}^{A_{m-1}}$ contains no copy of the sign-representation.
\end{proof}

\end{document}
